\section{Reading Guide}\label{sec:reading_guide}

Different readers have different backgrounds and goals. This guide suggests reading paths through the wiki depending on your needs.

\subsection{Path 1: Complete Beginner in Cosmology}\label{subsec:guide_beginner}

If you're new to cosmology and power spectra, follow this order:

\begin{enumerate}
    \item \cref{ch:background} (\nameref{ch:background}): Learn the expanding universe, distances, and growth of structure
    \item \cref{ch:power_spectrum}, Section ``Cartesian Fourier'': Power spectrum basics in flat space
    \item \cref{ch:power_spectrum}, Section ``Statistical Framework'': Ensemble statistics
    \item \cref{ch:power_spectrum}, Section ``Matter Perturbations'': What we're measuring
    \item \cref{ch:sfb} (\nameref{ch:sfb}): Spherical coordinates and basis functions
    \item \cref{ch:3d_ps} (\nameref{ch:3d_ps}): 3D power spectrum in spherical coordinates
\end{enumerate}

\textbf{Estimated reading time}: 2--3 weeks (assuming 30--60 min/day)

\subsection{Path 2: Experienced Cosmologist, New to SFB}\label{subsec:guide_experienced}

If you know power spectra but want to learn the SFB framework:

\begin{enumerate}
    \item Skim \cref{ch:background}: Refresh on distances ($\chi$, $D_A$, growth $D_+$)
    \item Skim \cref{ch:power_spectrum}: You know this material already
    \item \cref{ch:sfb}: Full reading—this is the new material
    \item \cref{ch:3d_ps}: Connect SFB to observables
    \item \cref{ch:observables}, Sections ``Angular Power Spectrum'' and ``SFB Kernels'': See how SFB predicts observations
\end{enumerate}

\textbf{Estimated reading time}: 3--5 days

\subsection{Path 3: Weak Lensing Focus}\label{subsec:guide_lensing}

If you're interested in weak lensing specifically:

\begin{enumerate}
    \item \cref{ch:background}, Section ``Comoving Distance'': Need distances for lensing kernels
    \item \cref{ch:sfb}, Sections on ``Spherical Harmonics'': Required for lensing fields
    \item \cref{ch:weak_lensing}: Complete reading
    \item \cref{ch:observables}, Sections on ``Redshift Binning'' and ``Cross-Spectra'': Practical observables
    \item \cref{ch:numerical}: Numerical integration for lensing kernels
\end{enumerate}

\textbf{Estimated reading time}: 4--5 days

\subsection{Path 4: Numerical Implementation}\label{subsec:guide_numerical}

If you're implementing code:

\begin{enumerate}
    \item \cref{ch:sfb}, Sections on ``Bessel Functions'': You'll code this
    \item \cref{ch:3d_ps}: Understand what you're computing
    \item \cref{ch:observables}, Section ``SFB Kernels'': The integrals you need to implement
    \item \cref{ch:numerical}: All of it—critical for your code
    \item Return to \cref{ch:weak_lensing} for kernel definitions (if including lensing)
\end{enumerate}

\textbf{Estimated reading time}: 1--2 weeks (+ coding time)

\subsection{Quick Reference for Specific Questions}\label{subsec:guide_quick}

\begin{itemize}
    \item \textbf{``What is the Limber approximation?''} → \cref{ch:3d_ps}, find ``Limber''
    \item \textbf{``How do redshift bins affect power spectra?''} → \cref{ch:observables}, Section ``Redshift Binning''
    \item \textbf{``What's the difference between convergence and shear?''} → \cref{ch:weak_lensing}, Section ``Convergence and Shear''
    \item \textbf{``How do I compute $C_\ell$ from $P(k)$?''} → \cref{ch:3d_ps}, find projection integral
    \item \textbf{``Which quadrature rule should I use?''} → \cref{ch:numerical}, Section ``Quadrature Strategies''
    \item \textbf{``What are E-modes and B-modes?''} → \cref{ch:weak_lensing}, Section ``Spin-Weight Harmonics''
\end{itemize}

Use the **Concept Index** (\nameref{sec:concept_index}) to find all relevant sections for a topic.

\subsection{Study Tips}\label{subsec:guide_tips}

\begin{enumerate}
    \item \textbf{Take notes}: Write equations in your own words; copy derivations by hand
    \item \textbf{Check dimensions}: Every equation should have consistent units—this catches errors
    \item \textbf{Follow cross-links}: When you see a reference like \cref{eq:sfb_expansion}, click to understand context
    \item \textbf{Work through examples}: If an example is given, compute it yourself before reading the answer
    \item \textbf{Question assumptions}: When a result is derived, ask ``why is this assumption valid?''
    \item \textbf{Connect to code}: If you're implementing, look at your code as you read—does it match the math?
\end{enumerate}

\subsection{Common Misconceptions (Clarified)}\label{subsec:guide_misconceptions}

\begin{description}
    \item[``SFB is just a change of coordinates''] Partially true, but SFB naturally incorporates spherical geometry and boundary conditions in a way Cartesian Fourier does not for spherical domains.
    
    \item[``3D power spectra are always larger than angular power spectra''] No—they contain different information. A 3D power spectrum can have more or less power in some modes depending on how information is distributed in redshift.
    
    \item[``Limber approximation is always accurate for large $\ell$''] Yes for $\ell \gg 1$, but the error depends on redshift distribution and wavenumber range. See \cref{ch:3d_ps} for limits.
    
    \item[``Weak lensing is only important for small angles''] Weak lensing is a full-sky effect. Small-angle approximations (Limber) are separately motivated and orthogonal to the lensing approximation.
    
    \item[``B-modes always indicate systematics''] Pure gravitational lensing produces only E-modes, but other physical effects (primordial gravitational waves, intrinsic alignments in some models) can produce B-modes. It's an observational diagnostic.
\end{description}

\subsection{Updating Your Knowledge as the Wiki Grows}\label{subsec:guide_updates}

This wiki is maintained continuously. When new sections are added:

\begin{enumerate}
    \item Check the Table of Contents for new chapters or sections
    \item Look for references to new material in existing sections (they'll be linked)
    \item Use the \nameref{sec:concept_index} to find related concepts
\end{enumerate}

---

\noindent
\textbf{Next}: See \nameref{sec:quick_ref} for a handy reference of key equations and formulas, then jump into the main content.
